\clearpage
\section{Build notes}
\label{sec:build}

\subsection{Canonical workspace}
From the repository root on Windows:
\begin{lstlisting}[language={}]
cmake -S . -B build -G "Visual Studio 17 2022" -A x64
cmake --build build --config Release
ctest --test-dir build -C Release -L IllumoWorkspace --output-on-failure
\end{lstlisting}
The build produces the static \texttt{Illumo} library,
\texttt{IllumoGameCore}, \texttt{IllumoGame}, \texttt{IllumoTests},
\texttt{IllumoGameTests}, and \texttt{IllumoPublicHeaderSmoke}. The aggregate
\texttt{IllumoRunTests} target refreshes both per-configuration test inventories
before running the workspace label. There is no compatibility executable named
\texttt{Illumo}.

The Python front end configures the root workspace and exposes the same flow:
Its interactive Windows console accepts left-clicks to activate rows, the left
setting arrow or right-clicks to cycle backward, and wheel input to navigate.
Mouse movement highlights the row under the pointer without activating it.
Keyboard controls remain available. Mouse hit testing requires the complete
dashboard to fit (at least 56 columns by 29 rows). Native console input mode is
restored around child commands and on exit; other terminals use keyboard input.
Build, test, coverage, tidy, and documentation menu actions display a live phase,
elapsed times, tool-reported progress, diagnostic line counts, and recent output.
Unknown totals use an activity indicator rather than an estimated percentage.
The complete combined log remains under \pathref{build-orchestrator-logs/}.
Ctrl+C stops the action's owned processes; Windows Job Objects retain descendant
ownership even if the launcher exits first. The completion screen shows exit
status and remains until Enter. Application launches and statistics retain
direct output, and explicit CLI output is unchanged.
\begin{lstlisting}[language={}]
python build.py build --config Debug --target IllumoGame --parallel
python build.py test
python build.py test --list-tests
python build.py test --test Illumo.Host.RequiredModuleRollback
python build.py test --test Illumo.SysCmdLine.WindowDimensionOptions
python build.py test --test IllumoGame.Config.CanvasCommandLineOptions
python build.py run -- -ww 1280 -wh 720
python build.py stats
python build.py stats --json
python build.py file-stats
python build.py file-stats -n 25
python build.py tidy
\end{lstlisting}
Exact runner cases execute in dedicated manual-testing directories so file-backed
configuration cannot mutate staged runtime data.

Named profiles are available for configure, build, test, run, and diagnostics:
\begin{lstlisting}[language={}]
python build.py profiles
python build.py doctor --profile release --json
python build.py profile-save dev --profile debug --no-docs --parallel 4
python build.py build --profile dev
\end{lstlisting}
Built-in debug and release profiles use separate build directories. Saved
settings live in Git-ignored \pathref{build-profiles.local.json}; an explicit
\texttt{--profiles-file} selects a shared version-1 JSON file. CLI values override
saved defaults, while extra CMake arguments append. Saving replaces only the
named entry, atomically, and excludes one-shot clean/fresh actions. The README
documents the schema and boolean overrides. Existing commands without a profile
keep their defaults; CMake remains authoritative.

The doctor command checks cache compatibility and tool availability with bounded
version probes, without configuring or repairing the tree. It does not prove
compiler/SDK readiness. Explicit generator and architecture conflicts fail before
configuration; foreign-source caches remain rejected even with fresh requested.
Command failures report elapsed time, working directory, and native exit status.
Orchestrator regression tests run through Python unittest discovery with
\texttt{-s tools -p test\_build.py}.

\subsubsection{Everyday development toolbox}
The compact dashboard opens a scrollable Development Tools submenu with a
test explorer, diagnostic browser, profile picker, existing-artifact shortcuts,
and source-file statistics. Arrows, wheel, Page Up/Down, and Home/End navigate;
mouse hit regions follow the visible viewport. Hover and ordinary item selection
only update details. Explicit action rows execute work. Slash focuses search;
ordinary characters (including navigation letters) become text, Backspace edits,
Enter applies, and Escape cancels. Outside text entry, Escape or q returns.

The explorer reads CTest's JSON inventory, including labels, commands, working
directories, and timeouts. Name search and project-label filters narrow the list.
Run Selected, Run Matching, Rerun Failed, and Refresh Inventory are explicit.
Preparation configures and builds declared discovery and smoke targets by
default; Run Existing skips preparation. Browsing never builds. Exact selected
configuration discovery files are mandatory: another-configuration fallback is
rejected, as is a mismatched single-configuration cache. Empty selections run
nothing. Escaped, anchored name filters are batched below Windows command-length
limits while CTest retains test working directories and timeouts.

Test runs use local \texttt{-T Test --no-compress-output} reporting without
submission. Fresh XML is saved per batch before a later batch can replace it.
Version-1 JSON sidecars beside raw logs record checkout, build directory,
configuration, effective settings, command directories, completion status,
durations, and available test results. Raw logs remain authoritative. Failed-test
reruns use only the latest test run with the same build identity. Incomplete or
cancelled results, missing tests, and legacy logs cannot silently supply reruns.

The diagnostic browser recognizes MSVC, Clang/clang-tidy, and CMake locations,
filters errors/warnings, and retains unrecognized output through raw log context.
Relative source paths resolve against recorded command directories. Previews are
read-only and explicitly labelled CURRENT SOURCE because files may have changed.
Missing files and invalid locations receive messages rather than guessed paths.

The profile picker shows built-in and saved effective settings. Applying a profile
carries its configuration, directory, generator, architecture, feature flags,
extra CMake arguments, and parallelism into ordinary dashboard commands. Manual
changes remain session overrides; switching profiles clears them. Saving requires
an explicit name and action. Coverage and tidy retain their dedicated Debug/Ninja
trees. Artifact shortcuts open only existing build directories, applicable logs,
coverage HTML, and generated PDFs through Windows default handlers; opening never
builds. All toolbox implementation remains standard-library-only in
\pathref{build.py} so independently generated projects retain the functionality.

The repository-statistics command reports the current Git branch, commit, and
working-tree counts together with tracked-file and categorized first-party line
counts. LOC means nonblank lines in the current tracked source, tests, shaders,
build tooling, documentation, and configuration files. Build trees,
\pathref{archive/}, \pathref{Illumo/thirdparty/}, \pathref{docs/output/}, and
binary assets are excluded. The interactive build console exposes the same
report, and \texttt{--json} provides a stable machine-readable form.
The \texttt{file-stats} command (or \texttt{stats --by-file}) provides a per-file
source breakdown sorted by LOC from largest to smallest to highlight large
components and god-object candidates, with optional \texttt{-n} count limits,
\texttt{--min-loc} thresholds, \texttt{--include-tests}, category filters, and JSON output.

\subsection{Standalone library}
The supported library-only build remains independently configurable:
\begin{lstlisting}[language={}]
cmake -S Illumo -B build-illumo
cmake --build build-illumo --config Release
ctest --test-dir build-illumo -C Release -L Illumo --output-on-failure
\end{lstlisting}
It builds the static library, \texttt{IllumoTests}, and the public-header
consumer smoke without seeing IllumoGame sources.

\subsection{Tests}
\texttt{IllumoTests} owns \texttt{Illumo.*}; \texttt{IllumoGameTests} owns
\texttt{IllumoGame.*}. Both support exact \texttt{--list} and \texttt{--run}.
CTest runs each case in an isolated working directory and selects a runner path
for the requested configuration in multi-config build trees. Headless tests do
not prove live OpenGL pixels, native dialogs, or native window lifecycle.

\subsection{Coverage (Clang/LLVM)}
\begin{lstlisting}[language={}]
cmake -S . -B build-coverage -G Ninja \
  -DCMAKE_BUILD_TYPE=Debug \
  -DCMAKE_C_COMPILER=clang -DCMAKE_CXX_COMPILER=clang++ \
  -DILLUMO_BUILD_DOCUMENTATION=OFF -DILLUMO_ENABLE_COVERAGE=ON
cmake --build build-coverage --target IllumoCoverage
\end{lstlisting}
The workspace target runs both inventories and public-header smoke, merges
profiles from both instrumented products, writes
\pathref{build-coverage/coverage-html/index.html}, and enforces the existing
85\% headless-testable first-party production-line gate. Tests, TestSupport,
vendored/system code, the concrete live window, and the OpenGL implementation
are excluded; proportional manual Windows smoke tests cover the live
graphics/dialog boundary.

\subsection{clang-tidy}
\begin{lstlisting}[language={}]
cmake -S . -B build-tidy -G Ninja \
  -DCMAKE_BUILD_TYPE=Debug \
  -DCMAKE_C_COMPILER=clang -DCMAKE_CXX_COMPILER=clang++ \
  -DILLUMO_BUILD_DOCUMENTATION=OFF -DILLUMO_ENABLE_CLANG_TIDY=ON
cmake --build build-tidy --target IllumoTidy
\end{lstlisting}
\texttt{ILLUMO\_ENABLE\_CLANG\_TIDY} defaults to ON, so first-party C++
compilation runs LLVM \texttt{clang-tidy} with \pathref{.clang-tidy} and treats
diagnostics as errors (D-T3). Disable with
\texttt{-DILLUMO\_ENABLE\_CLANG\_TIDY=OFF} or
\texttt{python build.py build --no-tidy}.
\texttt{python build.py tidy} configures a Ninja/Clang tree under
\texttt{build-workspace-tidy} and builds the batch \texttt{IllumoTidy} target
from \pathref{compile\_commands.json}. Vendored third-party sources are
excluded. \texttt{modernize-use-auto} stays disabled because house style
forbids \texttt{auto} (D-008).

\subsection{Runtime staging}
The Illumo-owned helper copies shaders, assets, notices, and licenses beside
\texttt{IllumoGame.exe}; tracked changes retrigger staging, and removed assets
are removed from the staged copy. IllumoGame's \texttt{envvars.json} uses
copy-if-missing so user runtime settings survive rebuilds.

\subsection{Documentation}
Run \texttt{python build.py docs} or \pathref{docs/build.ps1}. The script builds
\pathref{docs/output/illumo.pdf} and
\pathref{docs/output/architecture-map.pdf} from their LaTeX sources. The
optional \texttt{IllumoDocs} target uses the same script when TeX is available.
